% Generated draft: requires mathematical and source review.
\documentclass[11pt]{article}
\usepackage[margin=25mm]{geometry}
\usepackage{fontspec}
\setmainfont{Latin Modern Roman}
\usepackage{amsmath,amssymb,mathtools,unicode-math}
\setmathfont{Latin Modern Math}
\usepackage{enumitem,needspace,xurl,hyperref}
\setlength{\parindent}{0pt}
\setlength{\parskip}{5pt}
\setlength{\emergencystretch}{3em}
\newcommand{\N}{\mathbb{N}}
\newcommand{\Z}{\mathbb{Z}}
\newcommand{\Q}{\mathbb{Q}}
\newcommand{\R}{\mathbb{R}}
\newcommand{\C}{\mathbb{C}}
\newcommand{\F}{\mathbb{F}}
\newcommand{\A}{\mathbb{A}}
\newcommand{\PP}{\mathbb{P}}
\newcommand{\E}{\mathbb{E}}
\newcommand{\eps}{\varepsilon}
\newcommand{\dd}{\,\mathrm{d}}
\begin{document}
\section*{UnsolvedMath ID: 800987}
\textbf{Source problem code:} catalogue entry
\subsection{Definitions and mathematical statement}
Let \(S^{3}=\{x\in\R^{4}:\sum_{i=1}^{4}x_i^{2}=1\}\) be the standard smooth three-dimensional sphere. A Lorentzian metric on \(S^{3}\) is a smooth symmetric covariant \(2\)-tensor that is nondegenerate at every point and has one negative and two positive eigenvalues. The opposite overall sign convention may equivalently be used for both metrics.

Two metrics \(g\) and \(\bar g\) on the same manifold are \emph{projectively equivalent} if they have the same unparameterized geodesics: every nonconstant \(g\)-geodesic can locally be reparameterized by a smooth change of parameter with nonzero derivative to become a \(\bar g\)-geodesic, and conversely. This requirement includes geodesics of every causal type. The metrics are \emph{non-proportional} if there is no constant \(c\in\R\setminus\{0\}\) such that \(\bar g=cg\) on all of \(S^{3}\).

\textbf{Question (supplied source, Question 68).} Do there exist two smooth Lorentzian metrics \(g\) and \(\bar g\) on \(S^{3}\) that are projectively equivalent and non-proportional? Equivalently, is the following existence statement true?
\[
\exists\,g,\bar g\text{ Lorentzian metrics on }S^{3}
\quad\text{such that}\quad
\begin{cases}
 g\text{ and }\bar g\text{ have the same unparameterized geodesics},\\
 \bar g\ne cg\quad\text{for every constant }c\in\R\setminus\{0\}.
\end{cases}
\]
Here smoothness and the constant-multiple meaning of proportionality are conventional interpretations; the supplied excerpt does not explicitly specify their definitions. No additional geodesic completeness hypothesis is stated.

The source presents this question as a simplest version of the broader problem of understanding the topology of closed manifolds admitting non-proportional projectively equivalent metrics of indefinite signature. Its surrounding discussion reports local and global examples, with early examples attributed to Lagrange, Beltrami, and Dini; a local Riemannian description attributed to Levi-Civita; and a local arbitrary-signature description attributed to the reference labeled [Bol-Mat2015]. It describes the global Riemannian situation as comparatively well understood, citing [Mat-Top1998, Matveev2003a, Matveev2003b], while describing the global arbitrary-signature situation as largely unknown. These are supplied historical literature assessments, not independently verified claims of this draft.
\subsection{Short English statement}
Can the smooth three-dimensional sphere carry two Lorentzian metrics with exactly the same geodesic paths, allowing changes of parameter, such that neither metric is a constant multiple of the other? The supplied registry labels this question open; that label is not an independently verified current status.
\subsection{Sources}
\begin{itemize}
\item Local UnsolvedMath registry: supplied complete source record, permanent ID \(800987\), exact source problem code \texttt{catalogue entry}, title \emph{Open Problems and Questions about Geodesics --- Question 68}. The registry lists Burns and Matveev under attribution, gives proposed year 2019, and labels the record \texttt{open}. These are supplied metadata, not independently verified historical or mathematical conclusions.
\item Burns and Matveev, \emph{Open Problems and Questions about Geodesics}, identified in the supplied source list as 2019, Question 68: \url{https://arxiv.org/abs/1308.5417}. The supplied provenance records access on 2026-10-16 and extraction by \texttt{arxiv\_latex}; this is the recorded access date, not a claim that this draft independently accessed the source.
\item References appearing in the supplied excerpt: [Levi-Civita1896], credited there with the local Riemannian description; [Bol-Mat2015], credited with the local arbitrary-signature description; and [Mat-Top1998], [Matveev2003a], [Matveev2003b], cited for the global Riemannian situation. Their full bibliographic entries and individual URLs were not supplied. The excerpt also attributes early examples to Lagrange, Beltrami, and Dini without supplying individual citations.
\item Supplied upstream revision: \texttt{d371ee917ac9647fbb06cd6fc529deec6f3dda82}. The registry record is dated 2026-10-16. Its literature descriptions and open-status label are retained as supplied curation and have not been independently verified.
\end{itemize}
\end{document}
